\documentclass[12pt, oneside]{article}

\usepackage[T1]{fontenc}
\usepackage{palatino}
\usepackage{geometry}
\geometry{paper=letterpaper, top=1.25in, bottom=1.25in, left=1.25in, right=1.25in}

\usepackage{amsmath, amssymb, amsthm}
\usepackage{microtype}
\usepackage{setspace}
\usepackage{parskip}
\usepackage{titlesec}
\usepackage{hyperref}
\usepackage{array}
\usepackage{booktabs}
\usepackage{enumitem}
\usepackage{mdframed}

\hypersetup{colorlinks=true, linkcolor=black, citecolor=black, urlcolor=black}

\setstretch{1.15}

\titleformat{\section}{\large\bfseries}{}{0em}{}[\vspace{-0.3em}\rule{\linewidth}{0.4pt}]
\titleformat{\subsection}{\normalsize\bfseries\itshape}{}{0em}{}
\titlespacing{\section}{0pt}{2em}{0.8em}
\titlespacing{\subsection}{0pt}{1.2em}{0.4em}

\newtheorem{observation}{Observation}
\newtheorem{proposal}{Proposal}

\title{\textbf{Astrophysics After the Uniform Object}\\[0.5em]
\large Stratification, Burst Cycles, Reservoirs, and Reachability\\
from Compact Stars to the Cosmic Web}
\author{Flyxion\\[0.3em]\small Independent Researcher}
\date{June 2026}

\begin{document}

\maketitle
\thispagestyle{empty}

\begin{abstract}
\noindent
Across scales from nuclear-density matter to the cosmic web, observational astrophysics is converging on a structural shift: the uniform object --- a neutron star described by a single equation of state, a galaxy fixed in some spectral type, the cosmic web summarised by a two-point correlation function --- is giving way to layered, bursty, reservoir-constrained, and connectivity-structured descriptions that cannot be reduced to a single global parameter. This essay examines four concrete instances of that convergence, drawn from compact-star modelling~\cite{nouh2026}, high-redshift galaxy spectroscopy~\cite{roberts2026}, molecular gas detection at cosmic dawn~\cite{cescon2026}, and large-scale structure analysis~\cite{philcox2023}, and proposes that the RSVP, CLIO, Admissibility, and Frozen Processes frameworks~\cite{flyxion_frameworks} supply a structural vocabulary for the principle they collectively demonstrate. The central claim, read schematically rather than algebraically, is:
\[
\text{structure} = \text{constraint} + \text{reservoir} + \text{transition} + \text{reachable future}.
\]
More precisely, the essay argues that physical objects are projections of the intersection of admissible pasts and admissible futures, not fundamental isolated entities. The essay ends with open questions rather than equivalences.
\end{abstract}

\vspace{0.5em}
\noindent\rule{\linewidth}{0.4pt}
\vspace{0.5em}

\begin{center}
\begin{tabular}{>{\itshape}l l l}
\toprule
Paper & Astrophysical scale & Conceptual role foregrounded \\
\midrule
Nouh et al.~\cite{nouh2026} & stellar interior & admissible equilibrium \\
Roberts-Borsani et al.~\cite{roberts2026} & early galaxy evolution & bursty temporal phase \\
Cescon et al.~\cite{cescon2026} & ISM / fuel reservoir & latent reachability \\
Philcox \& Torquato~\cite{philcox2023} & cosmic web & percolating constraint geometry \\
\bottomrule
\end{tabular}
\end{center}

\smallskip
\noindent
The assignment of papers to terms is heuristic rather than exclusive. Each paper contains all four structural elements to varying degrees; the table identifies the aspect most strongly foregrounded by each analysis.

%-------------------------------------------------------------------
\section{Introduction: Astrophysics After the Uniform Object}
%-------------------------------------------------------------------

For most of its modern history, theoretical astrophysics has organised itself around a powerful simplification: the uniform object. A neutron star is a single thermodynamic phase characterised by one equation of state valid from centre to surface. A high-redshift galaxy is a particular spectral type occupying a fixed position in some classification scheme. The large-scale distribution of matter is a statistical field described by its two-point correlation function. The molecular content of a galaxy is a single scalar value derivable from an isolated tracer under standardised assumptions.

Stratified descriptions have existed in astrophysics for decades. Stellar structure models with separate core and envelope equations of state date from the mid-twentieth century; bursty star formation histories have been invoked since at least the 1990s; percolation ideas have appeared in cosmology for many years. What is new is not stratification as a concept but the observational precision that now makes the stratified picture unavoidable. JWST spectroscopy at $z \geq 10$, ALMA and VLA detections of CO in individual galaxies at $z > 7$, and large-scale dark matter simulations with $\sim 10^5$ halos make it no longer possible to suppress the heterogeneous structure of physical systems in favour of uniform-object approximations. This essay examines four cases where that suppression has become untenable --- compact-star modelling, high-redshift galaxy spectroscopy, molecular gas detection at cosmic dawn, and large-scale structure analysis --- and asks what structural principle they share.

Where older astrophysics explained systems as: \textit{object} $\to$ \textit{intrinsic properties} $\to$ \textit{observed behaviour}, the newer approach increasingly explains them as: \textit{configuration} $\to$ \textit{constraints} $\to$ \textit{trajectories} $\to$ \textit{observed behaviour}. A compact star is a region in parameter space carved out by stability conditions. A high-redshift galaxy is a point on a trajectory through burst space. A molecular gas reservoir is a state plus a quantified future capacity. The cosmic web is a connectivity manifold encoding the topology of the reachable structure of the universe. These are not more complicated descriptions of the same objects; they are a different kind of explanation.

The essay is offered as philosophical and mathematical analysis. The purpose is not to derive the astrophysical results from RSVP, CLIO, Admissibility, or Frozen Processes; rather, it is to show that the structural relationships uncovered by the four cases admit a common geometric interpretation when expressed in the language of those frameworks. The frameworks add an interpretive layer to work that does not need it to stand as science. The question is what structural grammar the four cases are collectively articulating, and whether naming that grammar illuminates something none of the individual analyses notices about itself.

%-------------------------------------------------------------------
\section{Compact Stars: Stratified Equilibrium and the CTOV Frontier}
%-------------------------------------------------------------------

\subsection{The framework}

Nouh et al.~\cite{nouh2026} replace a globally constant polytropic index with a radially varying one:
\begin{equation}
n(r) = \eta_0 - \eta_1 \tanh\!\left(\frac{r - r_c}{\varepsilon}\right),
\end{equation}
where $\eta_0 = (n_c + n_e)/2$, $\eta_1 = n_c - \eta_0$, $n_c$ is the core polytropic index, $n_e$ the envelope index, $\varepsilon$ the transition width, and $r_c$ the core boundary radius. The equation of state becomes $P = K\rho^{1+1/n(r)}$, and the resulting dimensionless CTOV system is solved via a Monte Carlo ensemble scheme that samples the boundary of reachable stellar configurations rather than tracing a single deterministic trajectory. The full parameter space is $\Theta = \{\rho_c, n_c, n_e, x_c, \varepsilon, \sigma\}$, where $\sigma = P_c/(\rho_c c^2)$ is the relativistic strength parameter. Table~11 of the paper calibrates this space against three nuclear equations of state: SLy4 ($n_c \approx 0.50$), APR ($n_c \approx 0.70$), and NL3 ($n_c \approx 1.0$), confirming that the parameter space is physically grounded.

\subsection{Stratification and the admissible mass-radius frontier}

RSVP does not treat stratification as an accidental complication layered on top of a uniform object; stratification is primary. The CTOV framework instantiates this precisely: observable stellar properties are determined more by the geometry of the transition region --- its width $\varepsilon$, its location $r_c$, the sharpness of the jump in $n(r)$ --- than by the bulk values of $n_c$ or $n_e$ in isolation. In RSVP language, $dn/dr$ at the transition boundary carries more explanatory information than $n$ far from it.

The paper's stability analysis identifies four conditions that jointly carve out the admissible region in parameter space: positivity and monotone decrease of density and pressure; adiabatic index $\Gamma > 4/3$ throughout; causality $v_s < c$; and the Zeldovich condition $0 < \sigma < 1$. For the soft-core models ($n_c = 0.5$, $n_e = 1$), the causality condition is violated at $\sigma > 0.3$; within the strictly causal range, the maximum reachable mass is only $M_{\text{max}} \approx 0.54$--$1.43\,M_\odot$ for $\varepsilon = 0.01$. The stiff-core models ($n_c = 1$, $n_e = 2$, $x_c = 0.7$) satisfy causality throughout and reach $M_{\text{max}} \geq 2.0\,M_\odot$ at $\sigma \geq 0.35$ for $\varepsilon = 0.01$, consistent with PSR~J0740+6620 ($2.08\,M_\odot$) and PSR~J0952--0607 ($2.35\,M_\odot$). Table~10 of the paper maps the reachable mass-radius frontier: configurations below the Oppenheimer--Volkoff floor ($M_{\text{max}} \approx 0.71\,M_\odot$) are unreachable by any stable compact star; those above the Rhoades--Ruffini causal ceiling ($M_{\text{max}} \leq 3.2\,M_\odot$) are unreachable by any causal star; the 21 observed neutron stars of Table~5 populate the admissible interior.

In Admissibility theory, the explanatory order is reversed: the admissible set is conceptually prior to any particular stellar configuration. The observable mass-radius distribution is the image, under an observational projection, of the admissible submanifold of parameter space. This is developed formally in Section~\ref{sec:math}.

\subsection{Parameters as compressed histories}

The Frozen Processes framework identifies the compact star as a compressed process: $\psi = C[\gamma]$, where $\gamma$ is the formation trajectory (collapse history, supernova, neutron star formation) and $C$ preserves only the dynamically stable information encoded in $(n_c, n_e, \varepsilon, x_c, \sigma)$. The core index $n_c$ is the preserved record of the nuclear interaction regime dominant during terminal compression; $n_e$ the preceding outer stratum; $\varepsilon$ the sharpness of the phase boundary; $x_c$ the fraction of stellar volume compressed into the inner stratum; $\sigma$ the relativistic depth of the final equilibrium. Reading these from observational mass-radius data is topological archaeology: reconstructing a formation trajectory from the compression $C[\gamma]$ it left behind.

%-------------------------------------------------------------------
\section{Early Galaxies: Bursty Time as a Hidden Coordinate}
%-------------------------------------------------------------------

\subsection{The framework}

Roberts-Borsani et al.~\cite{roberts2026} analyse 41 spectroscopically confirmed sources at $z \geq 10$ from JWST/NIRSpec prism observations, dividing the sample into C\,\textsc{iv}-strong sources (eight objects with detected C\,\textsc{iv}\,$\lambda\lambda 1548, 1550$\,\AA) and C\,\textsc{iv}-weak sources (33 objects). The C\,\textsc{iv}-strong sources occupy the extreme tails of nearly every observable distribution: the bluest UV continuum slopes ($\beta \lesssim -2.5$), the most compact morphologies ($r_{50} \lesssim 100$\,pc), and the highest star formation surface densities ($\Sigma_{\text{SFR}} \gtrsim 100\,M_\odot\,\text{yr}^{-1}\,\text{kpc}^{-2}$). Crucially they are not categorically separate: the two populations overlap continuously in the distribution tails.

The key quantitative result is the burstiness parameter $f_{\text{burst,3Myr}} = \text{SFR}_{0\text{--}3\,\text{Myr}}/\text{SFR}_{\text{tot}}$. For C\,\textsc{iv}-strong composites, $f_{\text{burst,3Myr}} \approx 0.69$; for C\,\textsc{iv}-weak composites, $f_{\text{burst,3Myr}} \approx 0.01$--$0.42$. The primary predictor of strong C\,\textsc{iv} emission is not metallicity (which varies by $\sim 0.4$\,dex), not dust content, and not AGN activity. It is the recency and intensity of star formation on timescales $\leq 3$\,Myr.

\subsection{Diversity as temporal stratification}

The paper's central argument is that spectroscopic diversity at $z \geq 10$ is driven by bursty star formation on short timescales. Strong C\,\textsc{iv} emitters are brief snapshots at the apex of a burst; sources devoid of emission lines are in the lull. Both are phases along a common evolutionary pathway.

This is temporal stratification. Just as the CTOV framework replaces a single global equation of state with a structured field of local equations of state varying in radius, Roberts-Borsani et al.\ replace a fixed galaxy type with a time-varying phase position along a burst/lull cycle. In RSVP terms, the burst and lull are two locally stable configurations of the coupled field $\mathcal{S} = (\Phi_{\rm SFR}, \mathbf{v}_{\rm gas}, S_{\rm ISM})$, and the gradient $d\Phi_{\rm SFR}/dt$ at the transition between them carries more explanatory information than the burst or lull values in isolation. The transition width --- the $\sim 3$\,Myr timescale over which C\,\textsc{iv}-emitting conditions are maintained --- plays the structural role of $\varepsilon$ in the CTOV framework, now acting in the temporal rather than the spatial coordinate domain.

\begin{observation}
The spectral diversity of $z \geq 10$ galaxies is a temporal cross-section through a single evolutionary trajectory. The C\,\textsc{iv}-weak/C\,\textsc{iv}-strong distinction is a phase-position indicator along that trajectory, not a taxonomic classification.
\end{observation}

In Frozen Processes language, the observed spectrum is a snapshot $\psi = C[\gamma_t]$ of the galaxy's trajectory at time $t$, not an intrinsic property. Two galaxies with identical present-day spectra might have radically different futures if one is at the beginning of a burst and the other at its end. The C\,\textsc{iv}-strong objects at $z \geq 10$ may be predecessors or successors of the C\,\textsc{iv}-weak objects, not categorical alternatives.

%-------------------------------------------------------------------
\section{Molecular Reservoirs: Latent Reachability in REBELS-25}
%-------------------------------------------------------------------

\subsection{The detection}

Cescon et al.~\cite{cescon2026} report the first detection of low-$J$ CO emission in an unlensed star-forming galaxy at $z > 7$: CO(3--2) at $3.4\sigma$ and CO(7--6) at $3.5\sigma$ in REBELS-25 ($z = 7.31$). The derived molecular gas mass is $M_{\text{mol}} = (1.0 \pm 0.4) \times 10^{11}\,M_\odot$ for $\alpha_{\text{CO}} = 3$, with the self-consistent \textsc{tuner} radiative transfer model yielding $M_{\text{mol}} = (1.8^{+1.0}_{-0.9}) \times 10^{11}\,M_\odot$ independently of $r_{31}$ and $\alpha_{\text{CO}}$. The galaxy is a dynamically cold rotating disc ($V_{\text{rot,max}}/\sigma = 11^{+6}_{-5}$) with near-solar metallicity ($Z \approx 0.85\,Z_\odot$), gas fraction $f_{\text{gas}} \approx 0.95$, and depletion timescale $\tau_{\text{dep}} \approx 1.2$\,Gyr.

\subsection{Reservoirs as reachable futures}

Approximately 95 per cent of the baryonic mass of REBELS-25 is cold molecular gas that has not yet formed stars. Define the reachability functional
\begin{equation}
\mathcal{R}(\psi) = \frac{M_{\rm mol}(\psi)}{M_\star(\psi)} = \frac{f_{\rm gas}}{1 - f_{\rm gas}}.
\end{equation}
For REBELS-25 with $f_{\rm gas} \approx 0.95$:
\begin{equation}
\mathcal{R} \approx \frac{0.95}{0.05} \approx 19.
\end{equation}
The galaxy contains roughly nineteen times more latent stellar mass than realised stellar mass. The depletion timescale $\tau_{\rm dep} \approx 1.2$\,Gyr specifies the rate at which the reachable future is being accessed. More precisely, $\mathcal{R}$ measures accessible future capacity under current dynamical conditions --- it is an upper bound on future stellar mass assembly rather than a guarantee, since feedback, gas ejection, and changing star formation efficiency all modulate the actual conversion. These are not descriptors of present state; they are coordinates on the reachable-future manifold of the system, subject to the constraints that keep the gas admissible as star-forming fuel.

\begin{observation}
The gas fraction $f_{\rm gas}$ and depletion timescale $\tau_{\rm dep}$ of a high-redshift galaxy are quantitative measures of its latent reachability: $\mathcal{R} = f_{\rm gas}/(1 - f_{\rm gas}) \approx 19$ bounds the ratio of accessible future stellar mass to present stellar mass under current dynamical conditions, and $\tau_{\rm dep}$ is the timescale over which that capacity is being accessed at the current star formation rate.
\end{observation}

\subsection{The reservoir as a compressed process record}

The near-solar metallicity ($Z \approx 0.85\,Z_\odot$) coexisting with $f_{\rm gas} \approx 0.95$ is in tension with simple closed-box evolution models. In Frozen Processes language this tension is informative: the present state $\psi = C[\gamma]$ is the compression of a formation trajectory $\gamma$ in which rapid enrichment and sustained cold accretion occurred concurrently. The gas reservoir encodes continuing inflow; the metallicity encodes chemical history; the rotating disc encodes accreted angular momentum; the dust mass encodes integrated stellar yields. The galaxy is not what it is. It is what its formation history has compressed into.

\subsection{The CMB as a redshift-stratified admissibility constraint}

At $z = 7.31$, $T_{\rm CMB} \approx 23$\,K is comparable to the excitation temperature of CO(3--2) ($T_{\rm ex} \approx 33$\,K). The CMB suppresses observable CO brightness by $f_{3\to 2} \approx 0.55$ and shifts the most populated rotational level from $J = 1$ to $J = 2$ in the best-fitting \textsc{tuner} model. The observational strategy admissible for tracing molecular gas is therefore stratified in redshift: CO(1--0) ceases to be the optimal probe at high enough $z$, and the boundary at which this occurs depends on the gas temperature and star formation activity. In Admissibility language, the set of viable observational projections $\pi_z : \mathcal{A} \to \mathcal{O}$ is itself a function of redshift --- an admissibility constraint imposed by the cosmological background on the map from admissible structure to observable data.

%-------------------------------------------------------------------
\section{The Cosmic Web: Disordered Heterogeneity and Percolation}
%-------------------------------------------------------------------

\subsection{The framework}

Philcox and Torquato~\cite{philcox2023} apply the mathematical formalism of disordered heterogeneous media to cosmological dark matter simulations, treating the galaxy distribution as a point cloud characterised by the pair correlation function $g_2(r)$, structure factor $S(k)$, void and particle nearest-neighbour functions $H_V(r)$ and $H_P(r)$, the pair-connectedness function $P_2(r;\eta)$, the direct-connectedness function $C_2(r;\eta)$, percolation thresholds, and the $\tau$ order metric.

On large scales ($r \gtrsim 200\,h^{-1}$\,Mpc), the galaxy distribution approaches hyperuniformity: $S(k) \sim k^{0.96}$ as $k \to 0$. On small scales it is almost antihyperuniform. The cosmological simulations percolate at $\eta_c = 0.252$, significantly below the Poisson threshold $\eta_c = 0.343$. Both systems share the fractal dimension $d_F \approx 2.4$ at the percolation threshold, implying the same universality class despite different physical origins. The pair-connectedness function $P_2(r;\eta)$ contains cosmological information equivalent to observing 25 times more galaxies when combined with $g_2(r)$, at minimal additional computational cost.

\subsection{The cosmic web as a connectivity manifold}

Using the formal definition from the paper, the pair-connectedness function is
\begin{equation}
P_2(r;\eta) \propto \mathbb{E}[\hat{\rho}(\mathbf{x})\hat{\rho}(\mathbf{x}+\mathbf{r})\,\hat{\Phi}(\mathbf{x},\mathbf{x}+\mathbf{r};\eta)],
\end{equation}
where $\hat{\Phi} = 1$ if there exists a path from $\mathbf{x}$ to $\mathbf{x}+\mathbf{r}$ remaining within the clustered phase, and $0$ otherwise. In CLIO terms, $P_2(r;\eta)$ can be read as a reachability statistic: the conditional probability that two galaxies are mutually accessible through the connected cosmic structure at reduced density $\eta$. The information advantage of $P_2$ over $g_2$ alone --- equivalent to 25 times more galaxies --- is the information advantage of knowing the reachability structure of the medium rather than only its correlation structure.

The percolation transition at $\eta_c$ marks the point at which global connectivity emerges from local clustering. In Admissibility language, one would say that the admissible configurations of the cosmic web simultaneously satisfy the large-scale hyperuniformity constraint imposed by primordial physics and the small-scale percolation structure imposed by gravitational clustering --- two projections of the same underlying constraint geometry onto different scale domains.

In RSVP terms, the gradients of the density field $\nabla\hat{\rho}$ at void boundaries and filament surfaces carry more information about the admissible structure of the cosmic web than the bulk values of $\hat{\rho}$ in clusters or voids. The lower percolation threshold of the cosmological simulations ($\eta_c = 0.252$ versus Poisson $\eta_c = 0.343$) is a signature of the enhanced gradient structure imprinted by primordial quasi-long-range correlations.

%-------------------------------------------------------------------
\section{Mathematical Derivations: From Parameters to Reachability}
\label{sec:math}
%-------------------------------------------------------------------

\subsection{Admissible configuration space, projection, and the pushforward measure}

Let $\Theta$ denote the parameter space of a physical system. For the CTOV case,
\[
\Theta_{\rm CTOV} = \{\rho_c, n_c, n_e, x_c, \varepsilon, \sigma\}.
\]
The admissible region is the subset satisfying all stability, causality, and thermodynamic constraints simultaneously:
\[
\mathcal{A} = \left\{\theta \in \Theta :
\rho(r) \geq 0,\;
P(r) \geq 0,\;
\Gamma(r) > \tfrac{4}{3},\;
v_s(r) < c,\;
0 < \sigma < 1
\right\}.
\]
The admissibility volume is
\[
V_{\mathcal{A}} = \int_{\mathcal{A}} d\mu,
\]
where $\mu$ is the natural measure on $\Theta$ induced by stellar formation processes. The mass-radius relation is an observational projection
\[
\pi_{\rm MR} : \mathcal{A} \to \mathbb{R}^2, \qquad \pi_{\rm MR}(\theta) = (M(\theta), R(\theta)).
\]
The observed neutron-star mass-radius distribution is the pushforward measure
\[
\mu_{\rm obs} = (\pi_{\rm MR})_*\,\mu_{\mathcal{A}},
\]
where $\mu_{\mathcal{A}} = \mu|_{\mathcal{A}}$ is $\mu$ restricted to the admissible submanifold. The observed distribution is therefore not a fundamental datum about objects; it is a projection of admissible structure, shaped by the geometry of $\mathcal{A}$ and the map $\pi_{\rm MR}$.

This formulation unifies the four papers' roles within the admissibility volume $V_{\mathcal{A}}$. The CTOV framework explores $\partial V_{\mathcal{A}}$ --- the boundary of the admissible region, mapped onto the mass-radius plane. Roberts-Borsani et al.\ explore motion within $V_{\mathcal{A}}$ --- the burst/lull trajectories are paths inside the admissible burst-space. Cescon et al.\ estimate the remaining volume of $V_{\mathcal{A}}$ accessible from the current state of REBELS-25, quantified by $\mathcal{R} \approx 19$. Philcox and Torquato study connectivity \textit{inside} $V_{\mathcal{A}}$ for the cosmic medium --- which configurations in the admissible space are mutually reachable through the percolating structure.

\subsection{Reservoir as reachable future}

The reachable stellar mass over a future interval $\Delta t$ satisfies
\[
M_\star(t + \Delta t) = M_\star(t) + \int_t^{t+\Delta t} \epsilon_{\rm SF}(t')\,M_{\rm mol}(t')\,dt',
\]
where $\epsilon_{\rm SF} = 1/\tau_{\rm dep}$ is the star formation efficiency per unit time. The reachability functional
\[
\mathcal{R}(\psi) = \frac{M_{\rm mol}(\psi)}{M_\star(\psi)} = \frac{f_{\rm gas}}{1 - f_{\rm gas}}
\]
is an observable: it can be computed directly from the CO-derived gas mass and the SED-fitted stellar mass. For REBELS-25 with $f_{\rm gas} \approx 0.95$, $\mathcal{R} \approx 19$. This number quantifies the scale of the reachable future relative to the present state; $\tau_{\rm dep} \approx 1.2$\,Gyr is the rate at which the reachable future is being accessed.

\subsection{Transition width as structural role}
\label{sec:transwidth}

Let a system pass between two locally coherent regimes $A$ and $B$ along a coordinate domain $\mathcal{D}$. A transition profile is a function $q : \mathcal{D} \to [0,1]$. Depending on the system, $\mathcal{D}$ may be spatial, temporal, or parameter-space valued:
\[
\mathcal{D} = \begin{cases}
\mathbb{R}_r, & \text{stellar radial stratification (Nouh et al.)}, \\
\mathbb{R}_t, & \text{galaxy burst/lull cycles (Roberts-Borsani et al.)}, \\
\mathbb{R}_\eta, & \text{percolation in density parameter space (Philcox \& Torquato)}.
\end{cases}
\]
A transition width is the characteristic scale $W$ over which $q$ changes substantially:
\[
W^{-1} \sim \max_{x \in \mathcal{D}}\!\left|\frac{dq}{dx}\right|.
\]
For the CTOV profile $q(r) = \tfrac{1}{2}[1 + \tanh((r - r_c)/\varepsilon)]$, this gives $W = \varepsilon$. The same structural role appears with $W = \tau_{\rm burst} \approx 3$\,Myr (Roberts-Borsani) and $W = \Delta\eta(L) \sim L^{-1/\nu}$ with $\nu \approx 0.85$ (Philcox--Torquato). These quantities have different dimensions --- $\varepsilon$ is a length, $\tau_{\rm burst}$ is a time, $\Delta\eta$ is dimensionless --- and should not be identified as the same parameter. They instantiate the same \textit{structural role}: the scale over which a system passes between two locally coherent regimes. Whether a unified field-theoretic account can derive all three from common first principles is an open question (Section~\ref{sec:open}).

\subsection{States as projected intersections of reachable pasts and futures}

For any physical state $\psi$, define the reachable past and future within $\mathcal{A}$:
\[
\mathcal{P}(\psi) = \{x \in \mathcal{A} : x \rightsquigarrow \psi\}, \qquad
\mathcal{F}(\psi) = \{y \in \mathcal{A} : \psi \rightsquigarrow y\}.
\]

\begin{proposal}[Reachability Decomposition]
The following should be understood as a conjectural organising principle rather than a derived result. Let $\mathcal{A}$ be the admissible configuration space of a physical system with observational projection $\pi : \mathcal{A} \to \mathcal{O}$. For any observed state $\psi \in \mathcal{O}$, the physically meaningful content of the state is determined by the structure of the triple $(\mathcal{P}(\psi),\; \psi,\; \mathcal{F}(\psi))$. Physical objects appear as projected intersections of reachable pasts and reachable futures:
\[
\psi = \pi\!\left(\overline{\mathcal{P}}(\psi) \cap \overline{\mathcal{F}}(\psi)\right),
\]
equivalently,
\[
\mathcal{A}_\psi = \overline{\mathcal{P}} \cap \overline{\mathcal{F}}, \qquad \psi = \pi(\mathcal{A}_\psi),
\]
where $\overline{\mathcal{P}}$ and $\overline{\mathcal{F}}$ denote closures within $\mathcal{A}$. The intersection --- rather than union or weighted measure --- is motivated by the requirement that a physically realised state must be simultaneously reachable from prior conditions and capable of further evolution: it must belong to both sets at once. This is a necessary condition, not a sufficient one, and the decomposition does not claim to be unique.
\end{proposal}

The four papers each supply one aspect of this structure. The CTOV framework supplies the reachable past: the formation history compressed into $(n_c, n_e, \varepsilon, x_c, \sigma)$ is readable from $\mu_{\rm obs} = (\pi_{\rm MR})_*\,\mu_{\mathcal{A}}$. Cescon et al.\ supply the reachable future: the molecular gas reservoir is quantified by $\mathcal{R} \approx 19$. Roberts-Borsani et al.\ supply the transition mechanism: the burst/lull cycle with timescale $\tau_{\rm burst} \approx 3$\,Myr is the process by which galaxies move through their admissible configuration space. Philcox and Torquato supply the connectivity geometry: $P_2(r;\eta)$ is the reachability measure of the cosmic admissible space at large scales.

%-------------------------------------------------------------------
\section{Synthesis: From Objects to Configuration Spaces}
%-------------------------------------------------------------------

\subsection{What each framework contributes}

RSVP contributes the ontological claim that observable structure is determined primarily by field gradients at transition boundaries rather than by bulk values within phases. The CTOV composite index $n(r)$ is a special case of an RSVP scalar field profile; the stellar properties that depend most sensitively on $\varepsilon$ are those determined by the gradient $dn/dr$ at the transition, not by $n_c$ or $n_e$ in isolation. This is a proposed extension beyond the CTOV paper --- RSVP would predict that $\varepsilon$ is not a free fitting parameter but is determined by the scalar sector of the underlying field theory, a claim the CTOV paper neither makes nor requires.

CLIO contributes the formal machinery of observational projection. The observable mass-radius distribution is the pushforward $\mu_{\rm obs} = (\pi_{\rm MR})_*\,\mu_{\mathcal{A}}$ of a formation measure on the admissible manifold. The reachability functional $\mathcal{R}$ is a coordinate on the reachable-future manifold. The pair-connectedness function $P_2(r;\eta)$ is a coordinate on the reachability structure of the cosmic admissible space. All three are CLIO projections of admissible structure onto observable subspaces. The four cases can be interpreted as probing different admissibility spaces --- $\mathcal{A}_{\rm CTOV}$, $\mathcal{A}_{\rm burst}$, $\mathcal{A}_{\rm mol}$, $\mathcal{A}_{\rm web}$ --- that instantiate a common geometric structure rather than constituting a single universal configuration space.

Admissibility contributes the reversal of explanatory order. The admissible sets $\mathcal{A}_{\rm CTOV}$, $\mathcal{A}_{\rm burst}$, and $\mathcal{A}_{\rm web}$ are conceptually prior to any particular neutron star, galaxy, or cosmic web realisation. The observable diversity in each case is a projection of admissible structure, not a catalogue of objects. Admissibility theory asks not ``what are the properties of this object?'' but ``what is the geometry of the admissible space this object inhabits, and where within that geometry does the object sit?''

Frozen Processes contributes the claim that $\psi = C[\gamma]$: the object is a compressed process, and the compression operator $C$ preserves only dynamically stable information from the formation trajectory $\gamma$. A neutron star is a collapse process compressed to its equilibrium parameters. REBELS-25 is an assembly and enrichment process compressed to $f_{\rm gas}$, $Z$, $V_{\rm rot}/\sigma$. A $z \geq 10$ galaxy in a lull is a burst process compressed to its current $f_{\rm burst,3Myr}$. The cosmic web is a gravitational collapse process compressed to its percolation threshold $\eta_c$ and pair-connectedness function $P_2$. In each case the compression $C$ is what observation performs: it projects the full dynamical trajectory onto the admissible manifold and returns the CLIO image $\pi(\mathcal{A}_\psi)$.

\subsection{The unifying geometric picture}

The four frameworks together propose that physical reality is best described not by a collection of objects with intrinsic properties but by a single geometric structure:
\[
\psi = \pi(\mathcal{A}_\psi), \qquad \mathcal{A}_\psi = \overline{\mathcal{P}} \cap \overline{\mathcal{F}} \subset \mathcal{A},
\]
where the field structure of RSVP determines the local geometry of $\mathcal{A}$, the compression of Frozen Processes determines $\overline{\mathcal{P}}$, the reachability functional $\mathcal{R}$ determines $\overline{\mathcal{F}}$, and the projection $\pi$ of CLIO determines what of $\mathcal{A}_\psi$ is observable.

The four cases exhibit a recurring geometric pattern that admits a common interpretation at four different scales. None of them articulates the pattern explicitly. The grammar appears to be present in the data prior to its explicit formulation.

%-------------------------------------------------------------------
\section{Open Problems and Future Directions}
\label{sec:open}
%-------------------------------------------------------------------

\textbf{The origin of the transition width.} In the CTOV framework, $\varepsilon$ is imposed phenomenologically. In Roberts-Borsani et al., $\tau_{\rm burst} \approx 3$\,Myr is inferred from SED modelling but not derived from first principles. In Philcox and Torquato, $\nu \approx 0.85$ is measured but not explained from the underlying physics. As developed in Section~\ref{sec:transwidth}, these three quantities instantiate the same structural role in different coordinate domains while having different dimensions and physical origins. Whether a unified field-theoretic account can derive all three from common first principles --- connecting $\varepsilon$ to the nuclear interaction Lagrangian, $\tau_{\rm burst}$ to the stellar population physics of massive stars, and $\nu$ to the statistics of primordial density fluctuations --- is an open question.

\textbf{Unified admissibility conditions.} The four stability conditions of Nouh et al.\ are applied independently. The burst/lull cycle is identified empirically. The percolation threshold is measured but not predicted from the pair correlation function alone. Whether all three can be understood as projections of a single admissibility functional onto different observable subspaces has testable consequences: the three boundary surfaces should meet at a common point in the parameter space of any theory that unifies them.

\textbf{Reservoir-to-burst connection.} REBELS-25 has $\mathcal{R} \approx 19$ and a dynamically cold rotating disc. Galaxies in Roberts-Borsani et al.\ at $z \geq 10$ are bursty, compact, and characterised by $f_{\rm burst,3Myr} \approx 0.69$ at peak. Whether these are different phases of the same evolutionary trajectory is a concrete question addressable by upcoming JWST and ngVLA observations. If so, REBELS-25 can be interpreted as sitting at large $\mathcal{R}$ with small $f_{\rm burst,3Myr}$ --- a pre-burst or inter-burst state whose future population of burst events is latent in its molecular reservoir.

\textbf{Percolation as a formation modulator.} The cosmological galaxy distribution percolates at $\eta_c = 0.252$ rather than the Poisson value $0.343$. Whether the star formation activity of galaxies is modulated by their position within the percolating cosmic structure --- whether galaxies in percolating filaments have systematically different $\mathcal{R}$ or $\tau_{\rm burst}$ than those in isolated halos --- connects the large-scale reachability structure of the cosmic web to the small-scale burst physics of Roberts-Borsani et al.

\textbf{Numerical implementation of the Reachability Decomposition.} The CTOV Monte Carlo scheme produces a sample from $\partial V_{\mathcal{A}}$ by sampling reachable stellar surfaces. Extending this to RSVP field equations --- replacing the TOV equation with field equations for $\mathcal{S} = (\Phi, \mathbf{v}, S)$ --- would provide a first numerical computation of $V_{\mathcal{A}}$ and $\mu_{\rm obs} = (\pi_{\rm MR})_*\,\mu_{\mathcal{A}}$ for a concrete physical system. The 21 neutron stars of Nouh et al.'s Table~5, the 41 galaxies of Roberts-Borsani et al.'s Table~1, and the QUIJOTE galaxy sample of Philcox and Torquato would serve as observational anchors at three different scales.

%-------------------------------------------------------------------
\section{Framework Glossary}
%-------------------------------------------------------------------

\begin{mdframed}[linewidth=0.4pt, innerleftmargin=10pt, innerrightmargin=10pt,
                  innertopmargin=8pt, innerbottommargin=8pt]
\noindent The essay draws on four interrelated frameworks~\cite{flyxion_frameworks}, each making a distinct formal commitment. They are used here as interpretive vocabulary; the astrophysical conclusions do not depend on accepting them.

\smallskip
\noindent\textit{RSVP} (Relativistic Scalar-Vector Plenum). Stratification is primary, not accidental. A physical object is a metastable configuration of a coupled scalar-vector-entropy field $\mathcal{S} = (\Phi, \mathbf{v}, S)$. Observable structure is determined more strongly by the gradients $\nabla\Phi$, $\nabla\mathbf{v}$, $\nabla S$ at transition regions than by bulk values within coherent phases. Transition boundaries carry more explanatory information than the apparent bulk.

\smallskip
\noindent\textit{CLIO}. The fundamental object is the admissible configuration space $\mathcal{A}$, not any particular state. Observations are projections $\pi : \mathcal{A} \to \mathcal{O}$ of that space onto the observable manifold. Physical objects are images of admissible structure under observational projection, not primary entities.

\smallskip
\noindent\textit{Admissibility}. The explanatory order is reversed relative to ordinary dynamical systems theory. Rather than deriving constraints from trajectories, trajectories are interpreted as paths through a pre-existing admissible region $\mathcal{A}$ whose geometry --- volume, curvature, bottlenecks, connectivity --- determines what evolutions are possible. The admissible set is conceptually prior to any particular trajectory.

\smallskip
\noindent\textit{Frozen Processes}. An object is a compressed process: $\psi = C[\gamma]$, where $\gamma$ is a historical trajectory and $C$ is a compression operator preserving only dynamically stable information. The object is not a thing that has a history; it is the history, stabilised at its current stage.
\end{mdframed}

%-------------------------------------------------------------------
\section{Conclusion: From Objects to Reachable Processes}
%-------------------------------------------------------------------

The four papers examined in this essay originate from different problems, employ different techniques, and address different scales of astrophysical organisation. None of them uses the language of admissibility, reachability, or stratification in the sense developed in the RSVP\/CLIO family of frameworks. Yet they converge on a common geometric picture.

A compact star is a projection of the admissible structure of its parameter space, $\mu_{\rm obs} = (\pi_{\rm MR})_*\,\mu_{\mathcal{A}}$, whose parameters encode the compression $C[\gamma]$ of the process that produced it. A high-redshift galaxy is a point on a trajectory through burst space, with spectral diversity reflecting temporal position rather than categorical distinction. A molecular gas reservoir is a quantified reachable future: $\mathcal{R} = f_{\rm gas}/(1-f_{\rm gas}) \approx 19$ for REBELS-25. The cosmic web is a percolating heterogeneous medium whose pair-connectedness function $P_2(r;\eta)$ maps the structural reachability of the cosmos at large scales.

\begin{observation}
A physical object is a projected intersection of its reachable past and its reachable future: $\psi = \pi\!\left(\overline{\mathcal{P}} \cap \overline{\mathcal{F}}\right)$. Its observable properties are determined by how much of its admissible past has been traversed, how much of its admissible future remains accessible, and what transitions connect the two. The grammar appears to be present in the data prior to its explicit formulation.
\end{observation}

\begin{thebibliography}{9}

\bibitem{nouh2026}
Nouh, M.\,I., El-Essawy, S.\,H., Foda, M.\,M., and Aboueisha, M.\,S.
\newblock ``Stochastic Analysis of Compact Stars under Composite Polytropes.''
\newblock \textit{Scientific Reports} \textbf{16}, 17858 (2026).
\newblock \href{https://doi.org/10.1038/s41598-026-53379-6}{doi:10.1038/s41598-026-53379-6}

\bibitem{roberts2026}
Roberts-Borsani, G., et al.
\newblock ``JWST Spectroscopic Insights into the Diversity of Galaxies in the First 500 Myr: Short-Lived Snapshots along a Common Evolutionary Pathway.''
\newblock \textit{Monthly Notices of the Royal Astronomical Society} \textbf{548}, 1--23 (2026).
\newblock \href{https://doi.org/10.1093/mnras/stag701}{doi:10.1093/mnras/stag701}

\bibitem{cescon2026}
Cescon, K., et al.
\newblock ``Direct Detection of Cool Molecular Gas in a Star-Forming Galaxy at $z = 7.31$.''
\newblock \textit{Monthly Notices of the Royal Astronomical Society} \textbf{549}, 1--20 (2026).
\newblock \href{https://doi.org/10.1093/mnras/stag924}{doi:10.1093/mnras/stag924}

\bibitem{philcox2023}
Philcox, O.\,H.\,E., and Torquato, S.
\newblock ``Disordered Heterogeneous Universe: Galaxy Distribution and Clustering across Length Scales.''
\newblock \textit{Physical Review X} \textbf{13}, 011038 (2023).
\newblock \href{https://doi.org/10.1103/PhysRevX.13.011038}{doi:10.1103/PhysRevX.13.011038}

\bibitem{flyxion_frameworks}
Flyxion.
\newblock Selected monographs developing the RSVP, CLIO, Admissibility, and Frozen Processes frameworks:

\medskip
\noindent\textit{Activation Manifolds as Admissibility Fields}.\\
\url{https://standardgalactic.github.io/playfloor/working/admissibility_fields.pdf}

\smallskip
\noindent\textit{Hidden Manifolds}.\\
\url{https://standardgalactic.github.io/alphabet/manuscripts/hidden_manifolds.pdf}

\smallskip
\noindent\textit{Constraint Before Content}.\\
\url{https://standardgalactic.github.io/playfloor/admissibility/constraint_before_content.pdf}

\smallskip
\noindent\textit{From Preservation to Reachability}.\\
\url{https://standardgalactic.github.io/animal-social/reachability.pdf}

\smallskip
\noindent\textit{Programs as Histories}.\\
\url{https://standardgalactic.github.io/playfloor/working/programs-as-histories.pdf}

\smallskip
\noindent\textit{Repair as a Fundamental Category}.\\
\url{https://standardgalactic.github.io/alphabet/research/repair_as_fundamental.pdf}

\smallskip
\noindent\textit{Frozen Processes}.\\
\url{https://standardgalactic.github.io/playfloor/reachability/frozen_processes.pdf}

\end{thebibliography}

\end{document}

